\section{Conclusion}\label{sec:conclusion}

The paper began after representation, with a map $F:X\to Y$. The missing architectural structure was not another formula but a marked receiver decomposition of the successor and a marked factorization of every receiver branch:
\[
\boxed{
B_j
=
\pi_j\chi_YF
=
G_jQ_j.
}
\]
Without the receiver marks, there are no separately addressed successor obligations. Without $Q_j$, receiver indexing with full shared access collapses to one product-valued map. Without $G_j$, the theory does not distinguish the information supplied to a receiver from the local computation performed after that information arrives. Bijective recodings of one interface are harmless re-presentations, while changes to its fibers alter which predecessor distinctions reach the receiver.

The factorization orders the later results. Interface fibers determine exact and approximate information barriers. Refining $Q_j$ into retained local state, source-resolved contributions, and aggregation distinguishes what is computed from what survives. In an additive linear regime, evaluated pair-contribution operators define a state-indexed operator family. Scalar relevance and transport become separate objects only when that factorization is itself marked.

Attention follows only after further restrictions. Pairwise scalar evidence, positive compositional linking, receiver-row ratio preservation, source-local payloads, additive aggregation, and shared finite separation rank yield masked query--key softmax. The supplement exhibits neighboring architectures obtained by removing these commitments one at a time. In the resulting positive scalar regime, effective-potential and free-energy expressions provide exact receiver-relative coordinates on an allocation that has already been derived; they are not the primitive definition of architecture.

The pointwise FFN enters on the other side of the central factorization. Its hidden coordinates are architectural because the specified local continuation computes a marked intermediate activation vector before output recombination. The FFN can construct nonlinear observables of the retained receiver state, but it cannot recover source distinctions removed before that state. A canonical PreNorm block is therefore, at a coarse grain, an attention-constructed receiver interface followed by a structured residual--FFN continuation. A finer grain exposes the intermediate residual state and the second receiver-wise update.

The resulting design space varies in receiver presentation, interface, contribution language, aggregation, local continuation, and schedule. Positive scalar routing adds a narrower set of allocation and potential coordinates. Variable successor presentations extend the same framework by allowing the marked target sector itself to depend on current state rather than treating a changed mask on one fixed carrier as creation of a new carrier.

Finally, an exact architectural presentation does not by itself establish an implementation boundary, an intervention-relative mediation relation, or identity across complete implementations. Those claims require additional evidence linking the marked presentation to implementation states, operations, admissible contrasts, invariance criteria, and tracking rules. The contribution here is the architectural object and the conditional reconstruction it supports: architecture comes before the familiar formula because the formula receives its architectural meaning from a marked organization of receiver access and continuation.
